% PPMI-Pipeline v2.1.8 poster source.
% Compile with pdfLaTeX, LuaLaTeX, or XeLaTeX.
% Required local theme files:
%   beamerthemegemini.sty
%   beamercolorthemegemini.sty
%   beamerinnerthemegemini.sty
% Required poster figures and QR assets belong in figures/.

\begin{filecontents*}{poster.bib}
@misc{ppmi_dataset,
  author       = {{Parkinson's Progression Markers Initiative}},
  title        = {{Parkinson's Progression Markers Initiative (PPMI) database}},
  year         = {2026},
  howpublished = {PPMI data portal},
  note         = {RRID:SCR\_006431}
}

@inproceedings{mckinney2010,
  author    = {McKinney, Wes},
  title     = {Data Structures for Statistical Computing in {Python}},
  booktitle = {Proceedings of the 9th {Python} in Science Conference},
  pages     = {56--61},
  year      = {2010},
  doi       = {10.25080/Majora-92bf1922-00a}
}

@article{pedregosa2011,
  author  = {Pedregosa, Fabian and Varoquaux, Ga\"{e}l and Gramfort, Alexandre
             and Michel, Vincent and Thirion, Bertrand and Grisel, Olivier
             and Blondel, Mathieu and Prettenhofer, Peter and Weiss, Ron
             and Dubourg, Vincent and Vanderplas, Jake and Passos, Alexandre
             and Cournapeau, David and Brucher, Matthieu and Perrot, Matthieu
             and Duchesnay, {\'{E}douard}},
  title   = {Scikit-learn: Machine Learning in {Python}},
  journal = {Journal of Machine Learning Research},
  volume  = {12},
  number  = {85},
  pages   = {2825--2830},
  year    = {2011}
}
\end{filecontents*}

\documentclass[final]{beamer}

\usepackage[T1]{fontenc}
\usepackage[utf8]{inputenc}
\usepackage{lmodern}
\usepackage[size=custom,width=120,height=72,scale=1.0]{beamerposter}
\usetheme{gemini}
\usecolortheme{gemini}

\usepackage{graphicx}
\usepackage{tikz}
\usetikzlibrary{arrows.meta,positioning}
\usepackage[numbers]{natbib}
\usepackage{url}

\graphicspath{{figures/}{./}}
\pdfstringdefDisableCommands{%
  \def\translate#1{#1}%
}

% Print-scale typography for a 120 cm x 72 cm conference poster.
\setbeamerfont{title in headline}{series=\bfseries,size=\fontsize{50}{56}}
\setbeamerfont{author in headline}{size=\fontsize{27}{32}}
\setbeamerfont{institute in headline}{size=\fontsize{21}{26}}
\setbeamerfont{block title}{series=\bfseries,size=\fontsize{30}{36}}
\setbeamerfont{block body}{size=\fontsize{24}{31}}
\setbeamerfont{caption}{size=\fontsize{15}{19}}
\setbeamercolor{architecture note}{fg=white,bg=geminiTeal}
\setbeamerfont{architecture note}{series=\bfseries,size=\fontsize{18}{23}}

% Three-column layout: (N + 1) * sepwidth + N * colwidth = paperwidth.
\newlength{\sepwidth}
\newlength{\colwidth}
\setlength{\sepwidth}{0.025\paperwidth}
\setlength{\colwidth}{0.3\paperwidth}
\newcommand{\separatorcolumn}{\begin{column}{\sepwidth}\end{column}}

\title{PPMI-Pipeline: A Reproducible Python Framework for Auditable Fixed-Horizon Longitudinal Modelling}
\author{Faris Haq}
\institute{Data Science \& Biomedical Research}

\begin{document}

\begin{frame}[t]
\begin{columns}[t]

\separatorcolumn

% -----------------------------------------------------------------------------
% Column 1: Objective and Methods
% -----------------------------------------------------------------------------
\begin{column}{\colwidth}

\begin{block}{Objective}
\begin{itemize}
  \setlength{\itemsep}{0.35em}
  \item \textbf{Software gap:} Longitudinal clinical analyses are often assembled as one-off scripts with implicit schemas, cohort rules, and temporal definitions that are difficult to audit or reproduce.
  \item \textbf{Computational risk:} Row-level evaluation can place observations from the same patient in both training and testing data, while undocumented preprocessing obscures the source of analytical outputs.
  \item \textbf{Primary contribution:} Deliver an open-source, configuration-driven Python architecture that makes ingestion, schema harmonisation, fixed-horizon construction, validation, modelling, and reporting explicit and testable.
  \item \textbf{Reference implementation:} PPMI clinical data provide the longitudinal use case \citep{ppmi_dataset}; the contribution is the reproducible computational method, not a clinical diagnostic or decision-support system.
\end{itemize}
\end{block}

\vspace{1cm}

\begin{block}{Methods}
\begin{itemize}
  \setlength{\itemsep}{0.35em}
  \item \textbf{Modular architecture:} \texttt{etl.py} validates and normalises source exports; \texttt{adapter.py} applies configuration-driven schema harmonisation for patient, visit, score, and date fields.
  \item \textbf{Fixed-horizon engine:} The earliest usable baseline (BL) is paired with the eligible follow-up nearest the configured horizon; baseline, target, score-change, and explicit exclusion records are generated.
  \item \textbf{Execution harness:} ElasticNet-regularised logistic regression implemented with scikit-learn \citep{pedregosa2011} uses \texttt{Baseline\_Score} as its sole predictor. It is a structural smoke test of end-to-end execution, not an optimized clinical predictor.
  \item \textbf{Validation by construction:} GroupKFold partitions by patient identifier, and every fold is programmatically audited for train/test patient overlap.
  \item \textbf{Configuration and provenance:} Horizon, tolerance, endpoint, label threshold, model settings, cohort counts, fold metrics, and generated artifacts are centrally configured and retained for audit.
\end{itemize}
\end{block}

\vspace{1cm}

\begin{block}{PPMI Data Acknowledgement}
{\fontsize{13.5}{17}\selectfont
Data used in the preparation of this article was obtained on 2026-07-15 from
the Parkinson's Progression Markers Initiative (PPMI) database
(\href{https://www.ppmi-info.org/access-data-specimens/download-data}{%
\nolinkurl{www.ppmi-info.org/access-data-specimens/download-data}}),
RRID:SCR\_006431. For up-to-date information on the study, visit
\href{https://www.ppmi-info.org}{\nolinkurl{www.ppmi-info.org}}.

\vspace{0.35cm}

PPMI -- a public-private partnership -- is funded by the Michael J. Fox
Foundation for Parkinson's Research, and funding partners; including AbbVie,
Alamar Biosciences, Aligning Science Across Parkinson's (ASAP), Arrowhead
Pharma, Arvinas, AskBio, BIAL, BioArctic, Biohaven, BlueRock Therapeutics,
Bristol Myers Squibb, Calico Labs, Capsida Biotherapeutics, Critical Path
Institute, DaCapo Brainscience, Denali, Edmond J. Safra Foundation, Eli Lilly,
Gain Therapeutics, GE Healthcare, Genentech, GSK, Insitro, Johnson \& Johnson
Innovative Medicine, Lundbeck, Merck, Neumora, Neuron23, Novartis, Olink,
Regeneron, Roche, Sanofi, Tenvie, UCB, Vanqua Bio, Voyager Therapeutics, The
Weston Family Foundation.
}
\end{block}

\end{column}

\separatorcolumn

% -----------------------------------------------------------------------------
% Column 2: Software architecture and verification evidence
% -----------------------------------------------------------------------------
\begin{column}{\colwidth}

\begin{block}{End-to-End Software Architecture}
The release exposes six explicit interfaces. Each stage consumes a documented
contract and emits an auditable artifact for the next stage.

\vspace{1.25cm}

\begin{center}
\begin{tikzpicture}[
  pipeline/.style={
    draw=geminiTeal,
    fill=geminiTeal!9,
    line width=1.4pt,
    rounded corners=8pt,
    minimum width=9.3cm,
    minimum height=3.7cm,
    text width=8.4cm,
    align=center,
    inner sep=0.35cm,
    font=\fontsize{17}{21}\selectfont
  },
  flow/.style={
    -{Latex[length=4.5mm,width=3.2mm]},
    draw=geminiTeal,
    line width=2.0pt,
    shorten >=4pt,
    shorten <=4pt
  }
]
  \node[pipeline] (etl) at (-10.8,0) {
    \textbf{1. ETL}\\
    Header validation\\
    Type normalization
  };
  \node[pipeline] (adapter) at (0,0) {
    \textbf{2. Schema Adapter}\\
    Configured field mapping\\
    Stable column contract
  };
  \node[pipeline] (horizon) at (10.8,0) {
    \textbf{3. Fixed-Horizon Builder}\\
    Baseline selection\\
    Follow-up and attrition audit
  };
  \node[pipeline] (validation) at (10.8,-7.0) {
    \textbf{4. Validation Harness}\\
    Patient-grouped folds\\
    Explicit overlap guard
  };
  \node[pipeline] (model) at (0,-7.0) {
    \textbf{5. Execution Model}\\
    ElasticNet smoke test\\
    Fold-local preprocessing
  };
  \node[pipeline] (reports) at (-10.8,-7.0) {
    \textbf{6. Reports / Provenance}\\
    Fold diagnostics\\
    Coefficients and run artifacts
  };

  \draw[flow] (etl.east) -- (adapter.west);
  \draw[flow] (adapter.east) -- (horizon.west);
  \draw[flow] (horizon.south) -- (validation.north);
  \draw[flow] (validation.west) -- (model.east);
  \draw[flow] (model.west) -- (reports.east);
\end{tikzpicture}
\end{center}

\vspace{0.45cm}

\begin{beamercolorbox}[
  wd=\linewidth,
  rounded=true,
  shadow=false,
  sep=1.3ex,
  center
]{architecture note}
  Reproducibility boundary:
  configuration $\rightarrow$ deterministic transformation $\rightarrow$
  patient-isolated validation $\rightarrow$ retained evidence
\end{beamercolorbox}
\end{block}

\vspace{1cm}

\begin{block}{Verification Evidence}
\begin{itemize}
  \setlength{\itemsep}{0.42em}
  \item \textbf{Safety under optimization:} Critical validation checks raise explicit runtime exceptions and remain active under \texttt{python -O}; no execution-path safety rule relies on \texttt{assert}.
  \item \textbf{Patient isolation:} Every GroupKFold partition invokes a train/test overlap guard before model fitting. The reference workflow completed all five folds with zero overlapping patient identifiers.
  \item \textbf{Installability:} \texttt{pyproject.toml} defines the package, Python compatibility, pinned runtime dependencies, development tooling, console entry point, and synthetic package fixture.
  \item \textbf{Automated verification:} Thirteen unit tests cover schema adaptation, fixed-horizon construction, target generation, model execution, patient isolation, report persistence, and visualisation.
  \item \textbf{Transparent model scope:} The single-feature ElasticNet model is a technical execution harness used to verify end-to-end orchestration and report generation. It is not an optimized predictor or clinical performance claim.
\end{itemize}

\vspace{0.35cm}

\begin{center}
{\fontsize{19}{24}\selectfont
\textbf{Auditable evidence chain}\\[0.22cm]
\texttt{schema} $\rightarrow$
\texttt{configuration} $\rightarrow$
\texttt{cohort counts} $\rightarrow$
\texttt{fold audit} $\rightarrow$
\texttt{reports}
}
\end{center}

\end{block}

\vspace{1cm}

\begin{block}{Reproducible Release Interface}
\begin{itemize}
  \setlength{\itemsep}{0.34em}
  \item \textbf{Install:} \texttt{pip install -r requirements-lock.txt}, followed by \texttt{pip install -e . --no-deps}.
  \item \textbf{Execute:} \texttt{ppmi-pipeline --data-path <export.csv>} invokes the canonical, configuration-validated workflow.
  \item \textbf{Verify:} \texttt{python -m unittest discover -s tests -v} exercises the public synthetic fixture without controlled data.
  \item \textbf{Inspect:} Logs, cohort summaries, fold audits, coefficients, and diagnostic figures provide a retained execution record.
\end{itemize}
\end{block}

\end{column}

\separatorcolumn

% -----------------------------------------------------------------------------
% Column 3: Validation, Conclusion, and References
% -----------------------------------------------------------------------------
\begin{column}{\colwidth}

\begin{block}{Validation}
\begin{itemize}
  \setlength{\itemsep}{0.35em}
  \item \textbf{Patient isolation:} GroupKFold groups by patient identifier, and an explicit runtime exception prevents repeated-patient leakage across train/test partitions.
  \item \textbf{Leakage audit:} Every fold is independently checked; the reference execution produced zero overlapping patient identifiers in all five folds.
  \item \textbf{Automated quality assurance:} All 13 unit tests passed across schema adaptation, fixed-horizon construction, target creation, modelling, validation, and visualisation.
  \item \textbf{Traceability:} Versioned configuration, run provenance, cohort counts, fold metrics, coefficients, and figure-generation scripts create an inspectable chain from input schema to output artifact.
\end{itemize}
\end{block}

\vspace{0.8cm}

\begin{block}{Conclusion}
\begin{itemize}
  \setlength{\itemsep}{0.35em}
  \item \textbf{Software contribution:} PPMI-Pipeline converts longitudinal data preparation and fixed-horizon modelling into a modular, documented, version-controlled computational protocol.
  \item \textbf{Engineering assurance:} Explicit schemas, deterministic cohort rules, patient-isolated validation, automated tests, and provenance records reduce hidden analytical failure modes.
  \item \textbf{Benchmark interpretation:} The single-predictor ElasticNet model demonstrates pipeline execution and leakage control only; it is not proposed for diagnosis, treatment, or clinical deployment.
  \item \textbf{Open science:} Shared code, synthetic fixtures, centralized configuration, audit evidence, and figure-generation scripts support inspection, reuse, and extension to future methodological studies.
\end{itemize}
\end{block}

\vspace{0.8cm}

\begin{block}{References}
\nocite{ppmi_dataset,mckinney2010,pedregosa2011}
{\fontsize{15.5}{19.5}\selectfont
\RaggedRight
\bibliographystyle{plainnat}
\bibliography{poster}
}
\end{block}

\vspace{0.8cm}

\begin{block}{Open Science \& Metadata}
\begin{center}
  \begin{minipage}[t]{0.43\linewidth}
    \centering
    \href{https://github.com/farishaq626-cloud/parkinsons_pipeline}{%
      \includegraphics[width=4.8cm]{github_qr.png}%
    }\\[0.25cm]
    \textbf{GitHub repository}
  \end{minipage}
  \hfill
  \begin{minipage}[t]{0.43\linewidth}
    \centering
    \href{https://doi.org/10.5281/zenodo.21225810}{%
      \includegraphics[width=4.8cm]{zenodo_qr.png}%
    }\\[0.25cm]
    \textbf{Zenodo archive}
  \end{minipage}

  \vspace{0.5cm}

  \textbf{PPMI-Pipeline v2.1.8}\\[0.25cm]
  {\fontsize{17}{21}\selectfont
    \href{https://github.com/farishaq626-cloud/parkinsons_pipeline}{%
      \texttt{github.com/farishaq626-cloud/parkinsons\_pipeline}%
    }\\
    \href{https://doi.org/10.5281/zenodo.21225810}{%
      DOI 10.5281/zenodo.21225810%
    }
  }
\end{center}
\end{block}

\end{column}

\separatorcolumn

\end{columns}
\end{frame}

\end{document}
